\documentclass[10pt]{article}
\usepackage[margin=20mm]{geometry}
\usepackage{amsmath,amssymb,mathtools,fontspec,unicode-math,xurl,hyperref}
\IfFontExistsTF{DejaVu Serif}{\setmainfont{DejaVu Serif}}{\setmainfont{Latin Modern Roman}}
\setmathfont{latinmodern-math.otf}
\setlength{\parindent}{0pt}\setlength{\parskip}{6pt}\sloppy
\title{Technical verification of
One-relator hierarchies}
\author{AI-assisted technical review}
\date{}
\begin{document}
\maketitle
Abstract: An explicit ascending HNN extension contradicts the finite-orbit and universal-periodicity assertions of Proposition 6.5 in the generality expressly admitted in Section 6. The weaker existence statement needed downstream is not refuted by this example, but requires a valid replacement proof. Three additional local corrections concern compactness in the normal-form theorem, regularity of cyclic covers, and the terminal complexity value. The latter corrections preserve the intended finite-complex applications. Conclusions are limited to arXiv version 3 and do not certify the unacquired journal text. Scientific adjudication remains pending.\par
\section*{Identity and scope}
The reviewed work is Marco Linton, One-relator hierarchies, associated with Duke Mathematical Journal DOI 10.1215/00127094-2024-0040 (https://doi.org/10.1215/00127094-2024-0040). The actual input is arXiv:2202.11324v3, dated 27 June 2024, 35 PDF pages. All page references use that PDF’s physical, one-based numbering, which agrees with its printed numbering. The repository lists v3 as its latest version and describes it as the version to appear in Duke; that metadata is not a textual comparison with the journal article.\par
The review traced the main implications through the hierarchy construction, normal forms, stability, and cyclic-stabilizer arguments. Their notes are incorporated with explicit limits. The primary review checked the disputed hypotheses, reconstructed the counterexamples, and examined the cited source-page images. No author code was executed. An independently written script checked the finite edge paths used in Problem 2 for \(m=1,\ldots,12\); the infinite family is proved analytically below.\par
\section*{Claims and positive checks}
Theorem 7.1, p. 31, concerns a specified one-relator hierarchy: a hyperbolic quasiconvex hierarchy, an acylindrical hierarchy, and a \(\mathbb Z\)-stable hierarchy without Baumslag–Solitar subgroups are equivalent, with virtual specialness and quasiconvexity of the indicated subcomplex images. It does not assert that every one-relator group without a Baumslag–Solitar subgroup is hyperbolic. Theorem 7.2 adds negative immersions, using Corollary 6.9 to obtain stability.\par
The reduced-rank estimate of Theorem 2.11, p. 8, can be organized consistently. Write \(A_0=\pi_1(A)\), \(B_0=\pi_1(B)\), \(E=A_0\cap B_0\). For \(C<A_0\), \(D<B_0\), choose representatives \(b\) for \(D\backslash B_0/E\), followed by representatives \(a\) for \((E\cap D^b)\backslash A_0/C\). Then \[C\cap D^{ba}=C\cap(E\cap D^b)^a.\] Successive applications of strong inertia give the required bound by \(\operatorname{rr}(D)\). Intersections outside \(B_0A_0\) are cyclic or trivial under the cited intersection theorem and contribute zero reduced rank. Connectedness of the Magnus-subgraph intersection has a real role in this identification.\par
For regular cyclic covers, the interval condition on translates of a \(\mathbb Z\)-domain gives the contractible fibers used in the horizontal-map argument. The hierarchy terminates after the correction in Problem 4. In the normal-form proof, the two thin-triangle cases provide a polynomial endpoint bound; a direct correction of its final subpath estimate appears in Problem 2.\par
The external Louder–Wilton paper was checked at the statement level: Theorem 1.2 relates negative immersions to the two-generator subgroup condition, Theorem A gives coherence, and Theorem B gives the co-Hopf and finite-conjugacy-class properties used in the descending-chain application. Those external proofs were not reconstructed. The final isomorphism and subgroup-hyperbolicity corollaries also retain their cited external dependencies.\par
\section*{Findings}
Problem 1 (Finite segment-orbit bounds do not imply Proposition 6.5).\par
\subsection*{Location and claim}
Proposition 6.5, p. 25, claims that if \(s_{\mathbb Z}(\psi)=\infty\), there are at most \(\operatorname{rr}(A)\) \(H\)-orbits of biinfinite geodesics through \(H\) all of whose finite subsets have noncyclic stabilizer. It further claims that every such line is preserved by a hyperbolic element. Its proof infers the finite-orbit statement from Lemma 6.4 by a pigeonhole argument.\par
Section 6, pp. 23–24, admits strongly inert subgroups of Magnus subgroups; Remark 6.2 explicitly includes a free base group with trivial relation, with the whole base group regarded as a Magnus subgroup. These stated allowances are essential to the example.\par
\subsection*{Counterexample with all hypotheses checked}
Let \[H=F(a,b),\quad A=H,\quad B=\langle a^2,b^2\rangle,\qquad
 \psi(a)=a^2,\quad\psi(b)=b^2,\] and \[G=\langle a,b,t\mid tat^{-1}=a^2,\ tbt^{-1}=b^2\rangle .\] A nonempty reduced word remains nonempty under this substitution, so \(\psi:A\to B\) is an isomorphism. Both edge groups are strongly inert: \(A=H\) trivially is, and \(B\) has rank two, a class explicitly listed on p. 6. For this particular \(B\), its folded core is two subdivided loops sharing their only vertex of valence greater than two; the fiber-product degree bound gives the same conclusion. The subgroup \(B\) is proper, since its abelianized image is \(2\mathbb Z^2\), and has infinite index: finite index \(d\) in \(F_2\) would give rank \(d+1=2\), hence \(d=1\), a contradiction.\par
For every \(n\), the recursively defined collection contains \[[\psi^n(H)]=[\langle a^{2^n},b^{2^n}\rangle].\] These groups have rank two. Thus \(s_{\mathbb Z}(\psi)=\infty\) and \(\operatorname{rr}(A)=1\).\par
In the Bass–Serre tree, a vertex \(gH\) has one edge in one direction and infinitely many in the other. Choose the direction in which its children are \(ghtH\), with \(h\) ranging over \(H/B\). Their stabilizers satisfy \[\operatorname{Stab}(ghtH)=ghBh^{-1}g^{-1}
       \subset gHg^{-1}=\operatorname{Stab}(gH).\] The child stabilizer is also a conjugate of \(H\), hence is noncyclic of rank two. The pointwise stabilizer of every finite directed segment is therefore its terminal vertex stabilizer and has rank two.\par
Every infinite forward ray from \(H\) extends uniquely backwards to a biinfinite directed line. Choosing just two children at each stage gives continuum many such lines. Every finite subset lies in a finite directed segment, and hence has the stipulated noncyclic stabilizer. Since \(H\) is countable, each \(H\)-orbit is countable. There are continuum many orbits, contradicting the asserted upper bound one. The countable group \(G\) also has at most countably many axes of hyperbolic elements, so most of these lines cannot be preserved by any hyperbolic element.\par
\subsection*{Impact and repair boundary}
The example is compatible with Lemma 6.4: orbit bounds at every finite length do not supply a single group element identifying entire infinite lines. Elements identifying longer prefixes need not stabilize as the prefix grows.\par
Theorem 6.8, pp. 26–27, uses Proposition 6.5 to obtain a periodic line and then a bounded-rank descending chain. Corollary 6.9 and the negative-immersions application use that theorem. The defective proposition is therefore a substantive proof dependency. Nevertheless, this example has the standard \(t\)-axis \[\cdots,t^{-1}H,H,tH,t^2H,\cdots .\] It does not disprove the weaker existential consequence actually needed in Theorem 6.8. Its total group is not asserted here to be a counterexample to the main one-relator theorem. A sufficient repair would establish that weaker existence statement in the required class, or provide another route to stability. Deleting the universal wording without a new proof does not establish the repair. The finite-\(s_{\mathbb Z}\) hypothesis of Theorem 6.14 excludes this example.\par
Problem 2 (Compactness is missing from the general normal-form theorem).\par
\subsection*{Location and construction}
Theorem 5.7, p. 19, assumes hyperbolicity of the vertex cover, acylindricity on the Bass–Serre tree, and prefix-closed relative quasigeodesic normal forms. Its proof invokes Theorem 3.4, p. 10, which additionally requires a compact total graph of spaces. The general setup does not impose this requirement.\par
For \(n\in\mathbb Z\), take two boundary lines with vertices \(A_n,B_n\), consecutive edges \(a_n,b_n\), and rungs \(r_n\) from \(A_n\) to \(B_{2n}\). Fill each pentagon \[A_n,A_{n+1},B_{2n+2},B_{2n+1},B_{2n},A_n.\] This is a contractible, locally finite strip \(X\). Give every edge length one. The function \(q(A_n)=n\), \(q(B_n)=n/2\) is 1-Lipschitz; the \(A\)-boundary is isometric to a line and every vertex is within two of it. Thus \(X^{(1)}\) is hyperbolic.\par
Take the edge space \(C\) to be an integer line with embeddings \(c_n\mapsto A_n\), \(c_n\mapsto B_n\), over a one-loop underlying graph. All vertex and edge groups are trivial, so the total group is \(\mathbb Z\), acting freely, and hence acylindrically, on its Bass–Serre line.\par
Use unique tree paths in \[T_A=A\cup\{r_n:n\in\mathbb Z\}\cup\{b_{2n}:n\in\mathbb Z\},
 \qquad T_B=B\cup\{r_n:n\in\mathbb Z\}\] as \(\eta_A,\eta_B\), and put \(\eta_X=\eta_A\). Each omitted edge has a tree detour of length four. Every subpath has length at most \(4d+5\), including endpoints inside edges, so these are uniformly 5-quasigeodesic normal forms. Tree paths are prefix-closed. Cancellation of their common tail gives the full relative-normal-form condition, since \(A\) and \(B\) are subtrees of their respective trees. Equivariance holds because \(\pi_1(X)=1\).\par
Index the strips in the total universal cover by \(j\in\mathbb Z\), with vertical edges \(A_n^{(j)}B_n^{(j+1)}\). The induced normal form from \(A_{2^m}^{(0)}\) to \(A_0^{(0)}\) stays in the zeroth strip and has length \(2^m\). In contrast, follow \[A_{2^{m-j}}^{(j)}
 \longrightarrow B_{2^{m-j}}^{(j+1)}
 \longrightarrow A_{2^{m-j-1}}^{(j+1)}
 \quad (0\le j<m).\] After \(2m\) steps this reaches \(A_1^{(m)}\). Cross to \(A_0^{(m)}\), then return using rungs and inverse vertical edges at index zero. This path has length \(4m+1\). No fixed constant can bound \(2^m\) by \(K(4m+1)+K\).\par
\subsection*{Correction and compact applications}
Adding compactness to Theorem 5.7 restores the missing premise of Theorem 3.4. The finite complexes and domains in Theorem 5.9, pp. 22–23, possess it, so the example does not refute the intended hierarchy application.\par
There is also a repairable estimate on p. 21: comparing \(f(d(u,v))+f(d(v,z))\) with \(f(d(u,z))\) has no justification for an arbitrary intermediate \(v\). But the normal-form identity already proved gives \(\eta(u,z)=\eta_X(u,v)*\mathrm{tail}\). For the corresponding subpath \(h'\), taking \(K\ge1\) yields \[|h'|\le K|\eta_X(u,v)|+K+|\mathrm{tail}|
       \le K|\eta(u,z)|+K
       \le K f(d(u,z))+K.\] This polynomial bound suffices for the invoked subexponential-detour theorem. Thus the displayed arithmetic error does not by itself defeat the compact case.\par
Problem 3 (The definition of a \(\mathbb Z\)-cover needs regularity).\par
Section 4, p. 10, defines a \(\mathbb Z\)-cover only by connectedness and deck group isomorphic to \(\mathbb Z\). Proposition 4.6, p. 13, claims a finite \(\mathbb Z\)-domain for every such cover of a finite CW-complex. This formal scope is too broad.\par
Let \(X\) be the three-petal rose and set \[G=F(a,b,c)=K*\langle c\rangle,\quad K=\langle a,b\rangle,
 \qquad L=\ker(K\to\mathbb Z),\ a\mapsto1,\ b\mapsto0 .\] For the connected cover corresponding to \(L<G\), malnormality of the free factor \(K\) and \(L\ne1\) imply \(N_G(L)=K\). Its deck group is therefore \(K/L\cong\mathbb Z\), as required. But the deck quotient is the intermediate \(K\)-cover of \(X\), a noncompact locally finite graph: the cosets \(Kc^n\) are distinct. A finite domain \(D\) satisfying \(\mathbb ZD=Y\) would surject onto that quotient and make it compact, a contradiction.\par
The covers actually chosen from epimorphisms onto \(\mathbb Z\) are regular. Inserting regular in the definition excludes the example and matches the proof’s use of regularity. This is a local scope correction, not a counterexample to the intended hierarchy.\par
Problem 4 (The terminal complexity is \((0,-1)\)).\par
The complexity definition and Proposition 4.12’s proof, p. 15, use \[c(X)=\bigl(|\lambda|/\deg(\lambda)-|X_\lambda^{(0)}|,\ -\chi(X)\bigr).\] For \(\langle a\mid a^m\rangle\), \(m\ge1\), the first coordinate is zero and \(\chi(X)=1\), so \(c(X)=(0,-1)\), although the proof calls the terminal value \((0,0)\). Conversely, \(\langle a,b\mid a^m\rangle\) has complexity \((0,0)\) and admits a cyclic cover through \(a\mapsto0,b\mapsto1\).\par
For a connected underlying graph of rank \(r\ge1\), the second coordinate is \(r-2\ge-1\); the first is a nonnegative integer. The lexicographic order on \(\mathbb N_0\times\{-1,0,1,\ldots\}\) is well founded. At \(r=1\) the group is finite cyclic and the complexity is \((0,-1)\). At \(r\ge2\), abelianization has rank at least \(r-1\), hence a quotient \(\mathbb Z\). Proposition 4.11’s strict decrease therefore still gives finite termination. This local arithmetic correction has a complete repair.\par
\section*{Assessment and limits}
Assessment: Major. Problem 1 disproves a substantive auxiliary proposition used on the route to the main negative-immersions result. Its impact is a required replacement argument, not an established counterexample to the main theorem. Problems 2–4 are local corrections; the compact, regular-cover hierarchy applications survive the specific counterexamples above.\par
The cited results on Magnus intersections, combinations of hyperbolic spaces, subexponential detours, virtual specialness, and the isomorphism algorithm were not independently reproved. The cyclic-stabilizer and Grushko arguments were traced, but every subsidiary combinatorial choice was not certified. These findings are not claimed exhaustive. A bounded search found no paper-specific correction; absence of a retrieved correction is not evidence that none exists. Only the acquired author version is assessed.\par
This is an initial technical review. The recorded scientific decision is `pending`; the human-review list is empty. It is not independent acceptance or a human journal recommendation.\par
9 M. Linton, One-relator hierarchies, arXiv:2202.11324v3, 27 June 2024, https://arxiv.org/abs/2202.11324v3. L. Louder and H. Wilton, Uniform negative immersions and the coherence of one-relator groups, Invent. Math. 236 (2024), 673–712, https://doi.org/10.1007/s00222-024-01246-4.\par
\end{document}
